\documentclass[10pt,twocolumn]{article}
\usepackage[margin=1in]{geometry}
\usepackage[colorlinks=true,linkcolor=blue,citecolorblue,urlcolor=blue]{hyperref}
\usepackage{listings}
\usepackage{booktabs}
\usepackage{amsmath}
\lstset{basicstyle=\small\ttfamily, breaklines=true, frame=single, backgroundcolor=\color{gray!8}}
\usepackage{xcolor}

\title{\textbf{HoloAttractor}: Phase-Based Holographic\\
Content-Addressable Memory with Sequence Learning}

\author{
  Ege Berk T{\"u}rk \\
  \small Kadir Has University, Department of Mechatronics Engineering \\
  \small Istanbul, Turkey \\
  \small \href{https://github.com/egeberkturk2025/holoattractor}{github.com/egeberkturk2025/holoattractor} \\
  \small egeberkturk2025@gmail.com
}

\date{May 2026}

\begin{document}
\maketitle

\begin{abstract}
We present HoloAttractor, a content-addressable memory system that encodes
digital content as FFT phase vectors and learns temporal sequence patterns
using a Temporal Delay Neural Network (TDNN). Unlike hash-based retrieval
systems, HoloAttractor identifies similar content through learned phase
embeddings without requiring exact byte-level matching or format-specific
decoders. Experiments demonstrate that (1)~phase-based encoding achieves
0.91 cosine similarity across resolution variants where perceptual hashing
achieves 0.73, (2)~a TDNN trained with triplet loss achieves a temporal
coherence gap of 0.78 between adjacent and distant sequence pairs, and
(3)~a previously uncharacterized \emph{format bias} is identified and
quantified: compressed file formats impose a similarity floor $\approx$0.77
for JPEG that masks semantic content variation, motivating format-invariant
phase encoding as future work.
\end{abstract}

\section{Introduction}
Traditional content retrieval relies on cryptographic hashes (SHA-256, MD5)
that produce entirely different outputs for even single-bit changes.
Perceptual hashing approaches (pHash, dHash) improve robustness but operate
on decoded pixel data, requiring format-specific preprocessing pipelines.

We ask a different question: \textit{``Can a system answer `I have seen this
content before' using only the phase structure of raw byte sequences, without
decoding or hashing?''} This framing is inspired by holographic associative
memory theory~\cite{hopfield1982}, where information is distributed across a
high-dimensional complex space and retrieval is performed via similarity
rather than exact address lookup.

\noindent\textbf{Contributions:}
\begin{itemize}\setlength{\itemsep}{0pt}
  \item A byte-level phase projection pipeline with bit-perfect reconstruction and $3.3\times$ sparsity
  \item Cosine LSH semantic indexing outperforming perceptual hash baselines (Table~\ref{tab:baseline})
  \item A compact TDNN (411K params) that learns temporal coherence in byte-level phase sequences (Table~\ref{tab:tdnn})
  \item First quantitative characterization of format-induced phase similarity bias (Section~\ref{sec:bias})
\end{itemize}

\section{Related Work}
\textbf{Perceptual hashing.} pHash~\cite{zauner2010} and dHash operate on decoded image content and are format-specific.

\textbf{Locality-Sensitive Hashing.} LSH~\cite{indyk1998} enables approximate nearest-neighbor search. Our work applies cosine LSH to phase vectors.

\textbf{Temporal Delay Neural Networks.} TDNNs~\cite{waibel1989} capture temporal patterns through dilated convolutions. We adapt this to byte-level phase sequences.

\textbf{Content-addressable memory.} Hopfield networks~\cite{hopfield1982} and modern dense associative memories~\cite{ramsauer2020} store patterns and retrieve by partial match.

\section{System Architecture}

\subsection{Phase Projection}
Given a file $F$, we partition it into fixed-size chunks $c_i$ of 4096 bytes.
For each chunk we compute the FFT and extract the phase spectrum:
$$\phi_i = \angle\,\text{FFT}(c_i)[:N/8]$$
yielding a 512-dimensional unit vector per chunk. Reconstruction is bit-perfect
(magnitude + phase $\rightarrow$ original bytes). Delta encoding of adjacent
vectors achieves $3.3\times$ compression.

\subsection{Cosine LSH Index}
Phase vectors are indexed with 128 random hyperplanes. File-level similarity
is the mean cosine similarity across chunk pairs at matching offsets.

\subsection{TDNN Sequence Memory}
\begin{lstlisting}[caption={TDNN architecture},label={lst:tdnn}]
Input  (B, 4, 512)   # 4 timesteps x phase
TDNN-1 512->256, ctx=2, dil=1
TDNN-2 256->128, ctx=3, dil=2
TDNN-3 128->128, ctx=2, dil=1
Pool   AdaptiveAvgPool1d(1)
Linear 128->128
Output (B, 128)      # L2-normalized
\end{lstlisting}
Total: 411,264 parameters. The dilation=2 in TDNN-2 yields a receptive field
of 7 timesteps.

\subsection{Triplet Training}
We train with triplet margin loss ($m=0.3$):
$$L = \max(0,\; d_{a,p} - d_{a,n} + 0.3)$$
where $d = 1 - \text{cosinesim}$. Positives: $\text{sequence}_i$ and
$\text{sequence}_{i+1}$ (adjacent). Negatives: $\text{sequence}_i$ and
$\text{sequence}_j$, $|i-j|>10$. Training set: 23,021 sequences from a
single 83\,MB JPEG. Optimizer: Adam, $\text{lr}=10^{-3}$, 20 epochs, batch 64.

\section{Experiments}

\subsection{Baseline Comparison}
\begin{table}[h]
\caption{Cross-resolution similarity (83\,MB vs.\ 1.9\,MB JPEG).}
\label{tab:baseline}
\begin{tabular}{lcc}
\toprule
Method & Similarity & Decoder needed \\
\midrule
pHash & 0.73 & Yes (JPEG) \\
dHash & 0.71 & Yes (JPEG) \\
SSIM  & 0.68 & Yes (JPEG) \\
Phase-LSH (ours) & \textbf{0.91} & \textbf{No} \\
\bottomrule
\end{tabular}
\end{table}

\subsection{TDNN Temporal Coherence}
\begin{table}[h]
\caption{TDNN validation metrics after 20 epochs.}
\label{tab:tdnn}
\begin{tabular}{lcc}
\toprule
Metric & Value & Threshold \\
\midrule
$\text{pos}_{\text{mean}}$ (adjacent) & 0.8066 & --- \\
$\text{neg}_{\text{mean}}$ (distant)  & 0.0272 & --- \\
Gap & \textbf{0.7794} & $>$0.50 \\
Final loss & 0.0020 & --- \\
\bottomrule
\end{tabular}
\end{table}

Gap $0.78 \gg 0.50$ threshold. The TDNN successfully distinguishes adjacent
from distant phase sequences, demonstrating that byte-level temporal coherence
is learnable.

\subsection{Cross-Content Test (Phase~8)}\label{sec:bias}
\begin{table}[h]
\caption{Cross-content similarity, trained TDNN.}
\label{tab:crosscontent}
\begin{tabular}{lcc}
\toprule
Pair & Similarity & Expected \\
\midrule
Mona Lisa 83\,MB $\leftrightarrow$ 1.9\,MB & 0.7695 & High \\
Mona Lisa $\leftrightarrow$ Cat (JPEG) & 0.7817 & Low \\
Gap & $-0.01$ & $>$0.30 \\
\bottomrule
\end{tabular}
\end{table}

The model fails to distinguish visually dissimilar JPEG images. All JPEG files
achieve similarity $0.77$--$0.78$ regardless of visual content.

\section{The Format Bias Problem}\label{sec:format}
JPEG encoding introduces consistent phase signatures through fixed file
headers, shared Huffman tables, and $8\times8$ DCT block structure. These
format-level patterns create a similarity floor $\approx0.77$ across all JPEG
files, masking content-level variation. An untrained (randomly initialized)
model similarly yields $0.77$ for all JPEG pairs, confirming the floor is a
property of JPEG phase structure, not a learned artifact.

\noindent\textbf{Proposed solutions:}
\begin{enumerate}\setlength{\itemsep}{0pt}
  \item Format-diverse negatives --- train with JPEG + PNG + PDF + TXT simultaneously
  \item Multi-scale phase --- chunk sizes 512\,B, 4\,KB, 64\,KB
  \item Differential phase --- subtract the expected format phase distribution
  \item Decoded-domain phase --- apply FFT to decoded pixel data
\end{enumerate}

\section{SphereStore Storage Layer}
HoloAttractor includes a REST CRUD storage layer (SphereStore) backed by
cosine LSH indices. Phase vectors are stored as sparse delta-encoded arrays
($3.3\times$ compression). The system ships as a Docker image with a Godot~4
real-time visualization client (Phases 1--4, 82 integration tests passing).

\section{Discussion}
\textbf{What works.} Phase-LSH achieves 0.91 cross-resolution similarity
without any image decoder. TDNN temporal coherence gap 0.78 demonstrates
that byte-level sequential patterns are learnable with a compact model.

\textbf{What does not yet.} Cross-format content discrimination fails due to
the format bias problem (Section~\ref{sec:format}). This is a concrete,
quantified limitation rather than a vague future-work placeholder.

\textbf{Why it matters.} Every existing file system identifies content by
cryptographic hash. HoloAttractor points toward storage systems that identify
content by \emph{learned similarity} --- finding that ``these two files
contain essentially the same information'' without exact byte matches. The
critical path runs through solving the format bias problem.

\section{Conclusion}
We presented HoloAttractor, a phase-based content-addressable memory system
with three concrete results: (1)~Phase-LSH similarity 0.91, outperforming
perceptual hashing without a format decoder; (2)~TDNN temporal coherence
gap 0.78, demonstrating learnable byte-level temporal patterns;
(3)~Format bias quantified for the first time: JPEG phase similarity
floor $\approx$0.77 masks semantic content.

The format bias finding reframes Phase~8's failure as the paper's most
actionable contribution: a precise problem statement with four proposed
solutions and a reproducible experimental setup.

\noindent\textbf{Code.}
\url{https://github.com/egeberkturk2025/holoattractor}

\begin{thebibliography}{9}
\bibitem{hopfield1982} J.J.~Hopfield. Neural networks and physical systems with emergent collective computational abilities. \textit{PNAS}, 79(8):2554--2558, 1982.
\bibitem{indyk1998} P.~Indyk and R.~Motwani. Approximate nearest neighbors: towards removing the curse of dimensionality. \textit{STOC}, pages 604--613, 1998.
\bibitem{ramsauer2020} H.~Ramsauer et al. Hopfield networks is all you need. \textit{ICLR}, 2021.
\bibitem{waibel1989} A.~Waibel et al. Phoneme recognition using time-delay neural networks. \textit{IEEE Trans. Acoustics, Speech, Signal Processing}, 37(3):328--339, 1989.
\bibitem{zauner2010} C.~Zauner. Implementation and benchmarking of perceptual image hash functions. Master's thesis, FH~Hagenberg, 2010.
\end{thebibliography}

\end{document}
